% Chapter 19 -- Microcontrollers and the road ahead
\chapter{Microcontrollers and the road ahead}\label{ch:mcu}

You built a CPU; a microcontroller is one shrink-wrapped. Same
fetch--decode--execute idea --- plus program flash, RAM, and every
peripheral anyone ever needed baked into one silicon slab that sells
for pocket change. This final chapter is a hand-off: the honest
MCU-versus-FPGA decision, the sensor board's natural MCU variant,
and the map to where the next books take over --- because cpu8 is
a beginning, not an end. This chapter asks for nothing but reading:
no sims, no board, just judgment.

\section{An MCU is the same idea, shrink-wrapped}\label{sec:shrink19}

\index{microcontroller}\index{MCU}
Our cpu8 was 302 cells: a program counter, an instruction ROM, a
register file, an ALU, an output port. An AVR ATmega328 --- the
Arduino UNO's brain --- is those same pieces plus flash you can
rewrite ten thousand times, 2\,KB of real RAM, three timers, an
ADC, UART, SPI, I\textsuperscript{2}C, and an interrupt controller,
all on one die.

Everything you met is there. The accumulator is a register file.
The UART's shift register is our ch16 FSM drawn the same way.
The timers are prescalers. The ADC is a comparator plus a SAR
state machine --- our ch18 ``threshold input,'' done properly.
Nothing is new in kind; everything is new in polish.

\begin{example}[Where the \$1 question is already settled]
An ATtiny10 costs about sixty cents at quantity
--- \emph{unverified, price-checked 2026-10-11}. For almost every
``blink something'' or ``read a sensor'' product, the MCU wins
before the argument starts: one chip, minutes of development,
power-down current in the nanoamps.
\end{example}

\begin{checkpoint}
If a \$0.60 part can do the job and is easier to program, the
\$20 FPGA needs a reason --- it isn't a flex, it's a tool with a
niche.
\end{checkpoint}

\section{The honest decision}\label{sec:decision19}

\index{MCU vs FPGA}
The difference isn't \emph{which is better} --- it's \emph{what the
hardware is}. An MCU runs your program one instruction at a time on
a fixed machine. An FPGA \emph{becomes} whatever machine you draw,
and a hundred copies of it if you want them. Pick by what the
problem demands, not by what's flashier.

\begin{figure}[htbp]
\centering
\input{figures/ch19/fig_decision}
\caption{The decision table: pick by the problem's shape, not the
part's reputation. Regenerate: \ttfamily
figures/ch19/fig\_decision.tex.}
\label{fig:decision19}
\end{figure}

\begin{example}[Three products, three picks]
\textbf{(a) LED strip controller} --- thousands of updates per second
is fine for an MCU's \texttt{delay}; pick the \$1 MCU. \textbf{(b)
Oscilloscope trigger} --- a 50\,ns edge can't wait for a fetch
cycle; pick the FPGA, hardware detects the edge in parallel.
\textbf{(c) Toy that beeps and remembers five buttons} --- the MCU
sleeps at a microamp and wakes on a keypress; the FPGA's static
power loses instantly. Two-thirds of hobby projects are
MCU-shaped; the FPGA's niche is precision and parallelism.
\end{example}

\section{The sensor board, with an MCU}\label{sec:variant19}

\index{sensor board, MCU variant}
ch18's board feeds FILT to an FPGA pin configured as a comparator
--- the honest-but-crude path. Its natural MCU variant wires FILT
to an ADC pin instead, and the whole design turns into firmware:

\begin{lstlisting}[caption={The sensor board's MCU variant, as an
Arduino sketch (shown, not run here --- the same sketch the docs
publish verbatim).},language=C]
void setup() { Serial.begin(9600); }
void loop() {
  int raw = analogRead(A0);            // 0..1023 across 0..Vcc
  float volts = raw * 5.0f / 1024.0f;
  float celsius = (volts - 0.5f) * 100.0f;  // TMP36's law
  Serial.println(celsius, 1);          // degC, one decimal
  delay(500);
}
\end{lstlisting}

Nine lines versus a Verilog module, an FSM, and a UART. That's the
point: for this problem, the MCU is unambiguously right. The FPGA
version's justification is different --- you build the machine
that does the same thing \emph{and} learn what the MCU hides.

\section{What Arduino hides}\label{sec:hides19}

\index{Arduino abstraction}
The sketch's \texttt{analogRead(A0)} is three real operations:
write the ADC's MUX register to select A0, start a conversion and
wait for its flag, read a result register. The library call
\emph{is} those registers --- the datasheet's ADC chapter is the
authoritative version, not the tutorial. Same for
\texttt{delay(500)}: it's a timer peripheral counting hardware
clocks --- our ch16 prescaler wearing a friendly name.

\begin{example}[Same debounce, two languages]
A button that bounces for $\sim$10\,ms: in C you poll a pin and
count milliseconds; in Verilog you clock a shift register and act
on a stable level.

\begin{lstlisting}[caption={Debounce both ways (shown, not run):
the C that waits, the Verilog that filters.},language=C]
// C on the MCU: sample, count, confirm
if (digitalRead(BTN) == level) {
  delay(12);                       // ride out the bounce
  if (digitalRead(BTN) == level)   // still there? real press
      act();
}
// Verilog on the FPGA: a 12ms filter is a shift register
always @(posedge clk)
  sr <= {sr[15:0], btn};           // 16 taps at ~1kHz
wire stable = (&sr) | (~|sr);      // all-ones or all-zeros
\end{lstlisting}

The C version spends CPU time \emph{waiting}; the Verilog version
adds sixteen flip-flops and never waits at all. Neither is wrong
--- they're the two cost models you'll choose between for the rest
of the book's road.
\end{example}

\begin{example}[Reading a datasheet's front page]
The ATmega328's page one offers: ``20\,MIPS at 20\,MHz, 32\,KB
flash, 2\,KB SRAM, 1\,KB EEPROM, 10-bit ADC, 6 PWM channels.''
Read honestly: fast enough for anything a sensor does, RAM so
small that \texttt{Serial.println} of a float spends tens of
bytes of it, one ADC multiplexed across six channels --- and
``20\,MHz'' is a \emph{ceiling}, not a default. The datasheet
answers what the tutorial doesn't: what's real versus what's
abstracted.
\end{example}

\section{Soft cores: the CPU inside the FPGA}\label{sec:softcore19}

\index{soft core}\index{PicoRV32}
There's a third option that makes the decision interesting: run a
real CPU \emph{on} the FPGA. A soft core is a processor written
in an HDL and dropped into fabric --- PicoRV32 is a complete
RV32I core that fits in roughly two thousand LUTs, small enough to
sit beside our cpu8's peripherals on the Tang Nano 9K. The payoff
is both-ways freedom: C code where convenience wins, hardware
accelerators where cycles matter, in one chip.

\begin{figure}[htbp]
\centering
\input{figures/ch19/fig_softcore}
\caption{PicoRV32 inside the FPGA: firmware runs on real RISC-V,
your Verilog peripherals attach beside it. Regenerate: \ttfamily
figures/ch19/fig\_softcore.tex.}
\label{fig:softcore19}
\end{figure}

\begin{example}[The integration, honestly scoped]
Dropping a soft core in is a short commit: instantiate the core,
attach BRAM holding the firmware, wire a UART. The full picture:

\begin{lstlisting}[caption={The soft-core shape (shown; a real
PicoRV32 build is Book 7's material --- its repo is cited).}]
picorv32 core (        // RV32I in ~2k LUTs
  .clk(clk27), .resetn(rst_n),
  .mem_valid(), .mem_ready(), .mem_addr(),
  .mem_wdata(), .mem_rdata(), .mem_wstrb());
// memory: firmware BRAM + UART-mapped registers
\end{lstlisting}

What this buys: gcc-compiled C running on silicon you shaped ---
an RV32I instead of our 10-opcode toy, but the same
fetch--decode--execute you can now see through.
\end{example}

\section{The road ahead}\label{sec:road19}

\index{learning map}
This book ends at the edge of three continuations. Book 7 takes
the machine you built to a real pipelined RISC-V --- hazards,
forwarding, caches, branch prediction: every shortcut our
two-state FSM didn't need. The analog craft continues in the
two classic texts: \emph{The Art of Electronics} for the deep
bench wisdom, \emph{High-Speed Digital Design} (H\&H) for the
signal-integrity world where wires stop being ideal. And the
daily drill is HDLBits --- short Verilog exercises that turn the
writing muscle into a habit.

\begin{figure}[htbp]
\centering
\input{figures/ch19/fig_map}
\caption{Where each path continues: Book 7 for the CPU, the texts
for analog, HDLBits for the drills. Regenerate: \ttfamily
figures/ch19/fig\_map.tex.}
\label{fig:map19}
\end{figure}

The arc this book ran: charge and voltage to gates and adders,
adders to an ALU, an ALU to a datapath, a datapath to a CPU that
ran Fibonacci on real silicon --- then breadboards, a real PCB
design, and now the honest choice of what to build next. The
capstone wasn't the UART frame; it was that none of it is magic
anymore. Machine the next books build is bigger, but the shape is
the one you've already drawn.

\section{Common mistakes}\label{sec:cm19}

\begin{tripwire}[``MCU equals FPGA'']
They answer different questions: \emph{what program runs} versus
\emph{what machine exists}. Saying ``the FPGA is more powerful''
like it's the upgrade misses that it solves a different problem
--- and pays for it in power, cost, and weeks of development.
\end{tripwire}
\begin{tripwire}[Arduino's abstractions as the truth]
\texttt{digitalWrite} hides direction registers, port writes, and
clock enables; \texttt{delay} hides a whole timer peripheral. Fine
for a toy --- wrong to assume the hardware is simpler than the
library makes it look. The datasheet is the truth; the library is
a convenience.
\end{tripwire}
\begin{tripwire}[Power assumptions imported from sketches]
An Arduino running a blink at 16\,MHz \emph{awake} draws tens of
milliamps; the same chip asleep draws microamps. Sketches never
sleep, so beginners think MCUs are power-hungry --- read the sleep
modes before judging battery life.
\end{tripwire}
\begin{tripwire}[Choosing the fashionable answer]
``FPGA because it's cooler'' is a project killer: the hobbyist
version of this mistake is a two-week build for a task a
sixty-cent chip solves in an afternoon. Pick the niche that fits,
not the one that impresses.
\end{tripwire}
\begin{tripwire}[Treating ``soft core'' as free]
PicoRV32 in the fabric is not ``an FPGA and an MCU for one
price''--- the LUTs, the BRAM for firmware, and the debug effort
are all spent. It's the right answer when you need both hardware
and convenience in one chip, not a freebie.
\end{tripwire}
\begin{tripwire}[Forgetting the sensor chain]
On a real product the MCU's job ends at the ADC register --- the
op-amps, filters, and protection circuits the earlier chapters
built still apply. Digital polish doesn't exempt the analog
front-end.
\end{tripwire}

\begin{exercise}[easy]
Name three peripherals on a typical MCU we never built --- and
which earlier chapter's circuit is hiding inside each.
\end{exercise}
\begin{exercise}[easy]
Why does \texttt{analogRead} count as a software convenience on
top of real hardware rather than a separate function?
\end{exercise}
\begin{exercise}[medium]
A project needs 96 LEDs updated at 1\,kHz and a button read every
200\,ms. MCU or FPGA, and which two rows of the table decide it?
\end{exercise}
\begin{exercise}[medium]
The sensor board's MCU variant has one real limitation the FPGA
version doesn't --- what is it, and does it matter for a room
thermometer?
\end{exercise}
\begin{exercise}[hard]
Pitch or reject: ``put a soft core on the sensor board so the MCU
and FPGA are one chip.'' State your criteria, then the answer.
\end{exercise}

\section{Summary}\label{sec:sum19}

\begin{center}
\begin{tabularx}{\linewidth}{@{}lX@{}}
\toprule
\textbf{Thing} & \textbf{The rule of thumb}\\
\midrule
MCU & the same fetch--decode--execute, plus every peripheral
shrink-wrapped\\
pick by & latency, parallelism, I\/O, cost, dev time, power ---
the problem's shape\\
Arduino & libraries are conveniences; the datasheet is the truth\\
soft core & real CPU inside fabric --- both-ways freedom, real
costs\\
the road & Book 7 for the real pipeline, texts for analog, HDLBits
for drills\\
the arc & charge $\to$ gates $\to$ ALU $\to$ CPU $\to$ board ---
none of it is magic anymore\\
\bottomrule
\end{tabularx}
\end{center}

\section*{Additional exercises}
\addcontentsline{toc}{section}{Additional exercises}

\begin{exercise}[easy]
Which three hardware pieces on an MCU does our cpu8 already have
in miniature?
\end{exercise}
\begin{exercise}[easy]
The honest reason most toys are MCUs: name it.
\end{exercise}
\begin{exercise}[easy]
What does a soft core consume that a hard MCU doesn't?
\end{exercise}
\begin{exercise}[medium]
A data-acquisition product needs 8 channels sampled truly
simultaneously at 1\,MHz --- which side of the table wins, and
which two rows are decisive?
\end{exercise}
\begin{exercise}[medium]
Explain why an MCU ``sleeps at microamps'' claim and a ``sketches
never sleep'' observation aren't a contradiction --- and which
peripheral makes sleep useful.
\end{exercise}
\begin{exercise}[medium]
Read any MCU datasheet's front page: list the three numbers that
set its real ceiling, and the marketing number to discount.
\end{exercise}
\begin{exercise}[hard]
Cost the two sensor-board builds end-to-end (MCU variant vs FPGA
variant) in parts AND developer time --- where's the honest break-
even?
\end{exercise}
\begin{exercise}[hard]
Design the ch18 board's next revision for an MCU: which nets
change, which peripheral does the header now target, and which
part can the board drop?
\end{exercise}
\begin{exercise}[medium]
A drone's motor-control loop must close at 20\,kHz with
sub-microsecond jitter, but the same product needs a rich app
stack for logging and radio. Argue the split: MCU, FPGA, or
both --- and which row of the decision table settles it?
\end{exercise}
\begin{exercise}[medium]
Your team proposes replacing a working MCU design with a
custom soft core ``for elegance''. List the three honest costs
they're under-counting, and the one case where the swap is
right anyway.
\end{exercise}

\section*{Hints and answers}
\addcontentsline{toc}{section}{Hints and answers}
\begin{description}
\item[Peripherals (easy, inline)] timer/PWM (our prescaler), ADC
(our threshold, done properly), UART (our ch16 TX/RX), I2C/SPI,
interrupt controller.
\item[analogRead truth (easy)] it writes real registers and waits
on a real flag --- the library names the mechanism, it doesn't
replace it.
\item[96-LED project (medium)] FPGA --- timing precision and
parallel outputs are the problem; cost and dev time still argue
MCU, but the spec's gap wins.
\item[MCU variant limit (medium)] the ADC's resolution and sample
rate cap the signal --- irrelevant for a room thermometer,
decisive for faster signals.
\item[Soft-core pitch (hard)] accept if you need one chip doing
both real-hardware timing and convenient firmware --- reject if
either side is unnecessary; the costs are LUTs, BRAM, and bring-up.
\item[Drone split (medium)] hard-real-time loop $\to$ FPGA (jitter
row), app stack $\to$ MCU (dev-speed row): both, partitioned ---
a classic real product answer, not a cop-out.
\item[Soft-core elegance swap (medium)] LUT and BRAM fabric
consumed, verification hours, and toolchain support; right only
when you already need the FPGA and the MCU is the scarcer part.
\item[Miniature pieces (easy)] program counter, register file,
ALU --- plus ROM and an I\/O port if we count peripherals.
\item[Why MCUs win toys (easy)] unit cost at volume plus dev
speed --- the decision table's last two rows.
\item[Soft-core cost (easy)] fabric LUTs, BRAM for firmware, and
your debug hours --- nothing is free.
\item[8-channel DAQ (medium)] FPGA --- latency and parallelism are
decisive; an MCU's ADC multiplexes one channel at a time.
\item[Sleep vs sketch (medium)] both true: the part can sleep
deeply, the sketches keep it awake --- the fix is sleep modes and
wake-on-interrupt, not a different chip.
\item[Datasheet ceiling (medium)] flash/RAM sizes and real clock
rate set limits; headline ``20 MIPS'' is the marketing number at
max clock.
\item[Build cost (hard)] MCU: \$3--8 parts + hours; FPGA: board
already owned + days of verification --- break-even is whether
you're building one unit or a product.
\item[MCU revision (hard)] header drops a pin (no FPGA threshold
input), targets ADC+GPIO; the FILT RC stays (analog front-end is
still real work); R2's bleed stays.
\end{description}

\section*{Further reading}
\addcontentsline{toc}{section}{Further reading}

\begin{itemize}
\item IBM Think, \emph{When to use FPGA vs
Microcontroller}.\newline
\url{ibm.com/think/topics/fpga-vs-microcontroller} --- the
comparison this chapter's decision table draws. Intro.
Verified 200, 2026-10-11.
\item Arduino Documentation, \emph{Blink} built-in
example.\newline
\url{docs.arduino.cc/built-in-examples/basics/Blink/} --- the
canonical first sketch, shown honestly as a sketch (not hardware).
Reference. Verified 200, 2026-10-11.
\item YosysHQ, \emph{PicoRV32 --- a size-optimized RISC-V
CPU}.\newline
\url{github.com/YosysHQ/picorv32} --- the soft-core named in this
chapter; real RV32I in $\sim$2k LUTs. Reference. Verified 200,
2026-10-11.
\item HDLBits.\newline
\url{hdlbits.01xz.net} --- the daily Verilog drill this chapter
recommends. Reference. Verified 200, 2026-10-11.
\item Wikipedia, \emph{Microcontroller}.\newline
\url{en.wikipedia.org/wiki/Microcontroller} --- what the parts
are. Reference. Verified 200, 2026-10-11.
\item Wikipedia, \emph{RISC-V}.\newline
\url{en.wikipedia.org/wiki/RISC-V} --- the ISA Book 7 teaches.
Reference. Verified 200, 2026-10-11.
\item Nandland, \emph{FPGA vs.\ Microcontroller vs.\
ASIC}.\newline
\url{nandland.com/lesson-2-fpga-vs-micro-vs-asic/} --- the
three-way comparison done right. [unverified: 403 this session;
prior verified access]
\end{itemize}

\begin{rulebox}[{Watch instead}]
\begin{enumerate}
\item \emph{FPGA vs Microcontroller vs Processor --- Choosing the
Right Hardware} --- Uplatz, 7:48,
\url{https://www.youtube.com/watch?v=i1ObzbkVfSA} --- the decision
in a short video. Matches \S\ref{sec:decision19}.
Opened and confirmed 2026-10-11.
\item \emph{What's an FPGA?} --- Charles Clayton, 1:25,
\url{https://www.youtube.com/watch?v=iHg0mmIg0UU} --- the
two-minute answer to the question this whole book builds.
Matches \S\ref{sec:shrink19}. Opened and confirmed 2026-10-11.
\end{enumerate}
\end{rulebox}

\chapternotes{MCU-vs-FPGA framing follows IBM's explainer plus
this book's own cost model; Arduino sketch and peripheral
abstractions per Arduino's docs; PicoRV32 per its repo. Part
counts and the ATtiny price are honest-marked and date-stamped.
No new sims or synth runs --- this chapter is the hand-off: the
tools it names were exercised earlier in the book or are Book 7's
material.}
